\begin{prob}
    For each problem below, state the largest theoretical lower bound that you can think
    of and justify (that is, your bound should be ``tight'').
    Provide justification for this lower bound. You do not need to find an algorithm
    that satisfies this lower bound.

    \emph{Example}: Given an array of size $n$ and a target $t$, determine the index of $t$ in the array.

    \emph{Example Solution}: $\Omega(n)$, because in the worst case any
    algorithm must look through all $n$ numbers to verify that the target is
    not one of them, taking $\Omega(n)$ time.

    \begin{subprobset}
        \begin{subprob}
            Given an array of $n$ numbers, find the second largest number.

            \begin{soln}
                $\Omega(n)$.

                Any algorithm must look through all $n$ numbers to be sure it didn't miss
                the second largest number, and so must take
                $\Omega(n)$ time. Algorithm: loop through the numbers, keeping
                track of the maximum and the next largest number seen.
            \end{soln}
        \end{subprob}

        \begin{subprob}
            Given an array of $n$ numbers, check to see if there are any duplicates.

            \begin{soln}
                $\Omega(n)$

                Again, we have to look through all of the numbers in the worst case;
                if we don't, we can't be sure that it isn't a duplicate of a number we
                have looked at.

                It turns out that the actual, \emph{tight} theoretical lower bound for this
                problem is $\Omega(n \log n)$, under some assumptions.\footnote{See:
                \url{https://en.wikipedia.org/wiki/Element_distinctness_problem}}
                But justifying that lower bound is not trivial, and this question asked
                you to state the best TLB that you could justify.

                If you said that the
                TLB is $\Omega(n \log n)$ because you though of sorting the list and
                checking for consecutive equal elements, this is good thinking -- but
                in order to claim that the TLB is therefore $\Omega(n \log n)$, you'd
                also need to justify the assumption that sorting (or something like it)
                is \emph{necessary} of any optimal algorithm. That's not an easy claim
                to justify!
            \end{soln}
        \end{subprob}

        \begin{subprob}
            Given an array of $n$ integers (with $n \geq 2$), check to see if there is
            a pair of elements that add to be an even number.

            \begin{soln}
                $\Omega(1)$

                This can be done in constant time by checking only the first 2 numbers.
                If the first two are both even, or both odd, then the answer is ``yes'',
                since the sum will then be even. If the first two numbers are of different
                polarity -- one even and the other odd -- then the answer is ``no'' if
                the size of the list is 2, otherwise it is ``yes'' since the third number
                in the array must be either even or odd; pairing it with one of the first
                two numbers with the same polarity gives an even sum.
            \end{soln}
        \end{subprob}

    \end{subprobset}

\end{prob}
